SIEGLE, DANIEL E, M.S.
``Cytochrome P450 Inhibitor Classification with Statistical Learning (2015)'' Directed by Dr. Weifan Zheng, 55pp.
\vspace{1cm}


This project compares methods for predicting cytochrome P450 (CYP) inhibition from compound structure. CYPs are a natural first target for \textit{in silico} models because adverse drug-drug interactions caused by CYP inhibition are a leading reason for rejecting drug candidates. I built and compared several non-linear, high-dimensional classification models using a large, publicly available high-throughput luminescence assay (PubChem AID 1851) against five CYP isozymes (1A2, 2C9, 2C19, 2D6 and 3A4). The methods compared are a Bayesian Binary QSAR method from the Molecular Operating Environment (MOE) and three standard machine learning methods implemented in Python: \textit{k}-nearest neighbors, random forests and support vector machines. All methods performed well. Support vector machines outperformed MOE Binary QSAR on all five isozymes, and \textit{k}-nearest neighbors and random forests performed comparably to it, showing that freely available software can match a method widely accepted in industry.
